Put \(N=N_m\), define
\[
 S:=\{(g_m(q),f(q)):q\in K_{\epsilon}\}\subset\mathbb R^{N+1},
 \qquad H:=\operatorname{conv}(S),
\]
and write a point of \(\mathbb R^{N+1}\) as \((x,t)\) with \(x\in\mathbb R^N\). The set \(K_{\epsilon}\) is compact: it is the intersection of the compact simplex \(\Delta_3\) with the two closed half-spaces \(p(q)\geq\epsilon\) and \(p(q)\leq1-\epsilon\). Every denominator in \(f\) is bounded below by \(\epsilon\) on this set, so \(f\) is continuous there; each Bernstein coordinate is a polynomial, so \(g_m\) is continuous. Therefore \(S\) is compact.

For every \(q\in K_{\epsilon}\), the multinomial identity gives
\[
 \mathbf 1^{\top}g_m(q)
 =\sum_{c\in\mathcal I_m}B_c^{(m)}(q)
 =\left(\sum_{a,y}q_{ay}\right)^m=1.
\tag{1}
\]
Thus \(S\) lies in the affine hyperplane \(\{(x,t):\mathbf 1^{\top}x=1\}\), which has dimension \(N\). We now give the required finite-support reduction directly. Suppose a finite convex representation of a point of \(H\) has \(k>N+1\) positive weights \(w_j\) on points \(s_j\in S\). The \(k\) points are affinely dependent, so there is a nonzero vector \(a=(a_j)_{j=1}^k\) satisfying
\[
 \sum_{j=1}^ka_j=0,
 \qquad
 \sum_{j=1}^ka_js_j=0.
\]
After replacing \(a\) by \(-a\) if necessary, some \(a_j>0\). Let
\[
 \rho:=\min_{j:a_j>0}\frac{w_j}{a_j}>0,
 \qquad w'_j:=w_j-\rho a_j.
\]
Then \(w'_j\geq0\) for all \(j\), at least one new weight is zero, and
\[
 \sum_jw'_j=1,
 \qquad
 \sum_jw'_js_j=\sum_jw_js_j.
\]
Iteration leaves at most \(N+1\) positive weights. Consequently every point of \(H\) is the image of some element of the compact set \(\Delta_N\times K_{\epsilon}^{N+1}\) under the continuous barycenter map
\[
 ((w_j)_{j=1}^{N+1},(q^j)_{j=1}^{N+1})
 \longmapsto
 \sum_{j=1}^{N+1}w_j(g_m(q^j),f(q^j)).
\]
Hence \(H\) is compact.

We next identify \(H\) with all attainable moment--target pairs. Every finite convex combination is the integral under its finite atomic probability measure, so \(H\) is contained in the set of attainable pairs. Conversely, let \(\nu\in\mathcal P(K_{\epsilon})\) and put
\[
 z:=\left(\int g_m\,d\nu,\int f\,d\nu\right).
\]
If \(z\notin H\), compactness gives a closest point \(h\in H\). For any \(s\in H\), convexity puts \(h+u(s-h)\in H\) for \(0\leq u\leq1\). The right derivative at \(u=0\) of the squared distance from \(z\) is nonnegative, hence
\[
 (z-h)^{\top}(s-h)\leq0
 \quad\text{for every }s\in H.
\tag{2}
\]
Apply (2) to \(s=(g_m(q),f(q))\) and integrate with respect to \(\nu\). The result is
\[
 (z-h)^{\top}(z-h)\leq0,
\]
contradicting \(z\neq h\). Therefore
\[
 H=
 \left\{\left(\int g_m\,d\nu,\int f\,d\nu\right):
 \nu\in\mathcal P(K_{\epsilon})\right\}.
\tag{3}
\]
The preceding affine-dependence argument proves at the same time that every pair in (3) has a representing measure with at most \(N+1\) atoms.

Fix any \(b\in\mathcal B_{m,\epsilon}\). By the definition of the Bernstein moment body, at least one probability measure has moment vector \(b\). By (3), the section
\[
 H_b:=\{t\in\mathbb R:(b,t)\in H\}
\]
is therefore nonempty. It is compact because \(H\) is compact, and it is convex because \(H\) is convex. Hence \(H_b\) is a closed interval. Equation (3) and the definitions of \(\mathcal F_{m,\epsilon}(b)\), \(L_m(b)\), and \(U_m(b)\) identify this interval equation by equation as
\[
 H_b
 =\left\{\int f\,d\nu:\nu\in\mathcal F_{m,\epsilon}(b)\right\}
 =[L_m(b),U_m(b)],
\tag{4}
\]
and compactness shows that both extrema are attained. By the converse construction in \(\text{lem:scm-moment-reduction}\), every representing measure in (4) generates a member of \(\mathfrak M_{m,\epsilon}\) with the observed count law \(b\) and causal ATE equal to its objective. Thus (4) is the sharp causal fiber, not a relaxation, and the \(N+1\)-atom reduction gives compatible SCM witnesses for every pair and both endpoints.

It remains to prove the sharper dual conclusions when
\(b\in\operatorname{relint}(\mathcal B_{m,\epsilon})\). Let
\[
 A:=\operatorname{aff}(\mathcal B_{m,\epsilon}),
 \qquad V:=\operatorname{lin}(A-A),
\]
and work in the Euclidean affine space \(A\times\mathbb R\), whose translation space is \(V\times\mathbb R\). Define the upward closure
\[
 E_-:=H+\{(0,s):s\geq0\}.
\]
This set is convex. It is closed: if \(h_k+(0,s_k)\) converges with \(h_k\in H\) and \(s_k\geq0\), compactness supplies a subsequence \(h_k\to h\in H\), after which \(s_k\) also converges to a nonnegative limit. By (4),
\[
 E_- =\{(x,t):x\in\mathcal B_{m,\epsilon},\ t\geq L_m(x)\}.
\tag{5}
\]
For \(k\geq1\), set \(z_k=(b,L_m(b)-k^{-1})\notin E_-\), and let \(h_k\) minimize the distance from \(z_k\) to \(E_-\). Such a minimizer exists because \(E_-\) is closed and the search may be restricted to a bounded ball once any comparison point is fixed. Since \((b,L_m(b))\in E_-\),
\[
 0<\lVert h_k-z_k\rVert\leq k^{-1},
\]
so \(h_k\to(b,L_m(b))\). Put
\[
 n_k:=\frac{h_k-z_k}{\lVert h_k-z_k\rVert}\in V\times\mathbb R.
\]
For \(y\in E_-\), the segment from \(h_k\) to \(y\) lies in \(E_-\). Differentiating squared distance at the minimizing endpoint as above gives
\[
 n_k^{\top}(y-h_k)\geq0.
\]
The unit sphere is compact, so along a subsequence \(n_k\to(a,\beta)\in V\times\mathbb R\) with \(\lVert(a,\beta)\rVert=1\). Passing to the limit gives
\[
 a^{\top}(x-b)+\beta\{t-L_m(b)\}\geq0
 \quad\text{for every }(x,t)\in E_-.
\tag{6}
\]
The upward closure in (5) forces \(\beta\geq0\). If \(\beta=0\), then (6) says \(a^{\top}(x-b)\geq0\) for every \(x\in\mathcal B_{m,\epsilon}\). Because \(b\) is in the relative interior, a relative ball about \(b\) in \(A\) lies in the moment body. Applying the inequality to \(b\pm\delta v\) for every sufficiently small \(\delta>0\) and every unit \(v\in V\) gives \(a^{\top}v=0\). Thus \(a=0\), contradicting \(\lVert(a,\beta)\rVert=1\). Therefore \(\beta>0\), and division by \(\beta\) lets us take \(\beta=1\).

Define \(\lambda^-\in\mathbb R^N\) by
\[
 \lambda^-:=-a+\{L_m(b)+a^{\top}b\}\mathbf 1.
\]
Using (1), inequality (6) at \((g_m(q),f(q))\in H\subseteq E_-\) becomes
\[
 (\lambda^-)^{\top}g_m(q)
 =L_m(b)-a^{\top}\{g_m(q)-b\}
 \leq f(q)
 \quad(q\in K_{\epsilon}),
\tag{7}
\]
while (1) also gives
\[
 (\lambda^-)^{\top}b=L_m(b).
\tag{8}
\]
Thus \(\lambda^-\) is a feasible Bernstein minorant attaining the primal lower endpoint. Conversely, any feasible minorant \(\eta\) and any \(\nu\in\mathcal F_{m,\epsilon}(b)\) satisfy
\[
 \eta^{\top}b
 =\int\eta^{\top}g_m\,d\nu
 \leq\int f\,d\nu.
\]
Taking the infimum over \(\nu\) proves weak duality \(\eta^{\top}b\leq L_m(b)\), so (8) is the attained lower dual optimum.

For completeness, the upper argument is the sign-reversed construction rather than an appeal to symmetry. Let
\[
 E_+:=H-\{(0,s):s\geq0\}
 =\{(x,t):x\in\mathcal B_{m,\epsilon},\ t\leq U_m(x)\}.
\]
Project \((b,U_m(b)+k^{-1})\) onto this closed convex set and pass to a normalized limiting normal. Downward closure makes its last coefficient negative, while the same relative-interior argument rules out a zero last coefficient. After scaling, the resulting support inequality has the form
\[
 a_+^{\top}(x-b)-\{t-U_m(b)\}\geq0
 \quad((x,t)\in E_+).
\]
Consequently
\[
 \lambda^+:=a_++\{U_m(b)-a_+^{\top}b\}\mathbf 1
\]
satisfies
\[
 (\lambda^+)^{\top}g_m(q)
 =U_m(b)+a_+^{\top}\{g_m(q)-b\}
 \geq f(q),
 \qquad
 (\lambda^+)^{\top}b=U_m(b).
\tag{9}
\]
Integrating any feasible majorant against any measure in the fiber gives the reverse weak-duality inequality, so (9) is the attained upper dual optimum. Equations (7)--(9) prove
\[
 D_m(b)=(L_m(b),U_m(b))
\]
with no duality gap and with both extrema attained.

Finally, consider the lower optimum and its compact contact set
\[
 T_-:=\{q\in K_{\epsilon}:(\lambda^-)^{\top}g_m(q)=f(q)\}.
\]
Let \(\nu_-\in\mathcal F_{m,\epsilon}(b)\) attain \(L_m(b)\), whose existence was proved in (4). Equations (7) and (8) give
\[
 0=\int_{K_{\epsilon}}
 \{f(q)-(\lambda^-)^{\top}g_m(q)\}\,d\nu_-(q).
\]
The integrand is continuous and nonnegative, so it vanishes \(\nu_-\)-almost surely; hence \(\nu_-(T_-)=1\). The barycenter argument used for (3), now on the compact set \(g_m(T_-)\), gives
\[
 b\in\operatorname{conv}g_m(T_-).
\]
If \(r=\dim\operatorname{aff}g_m(T_-)\), the affine-dependence elimination already proved reduces a convex representation of \(b\) to at most \(r+1\) contact points. The resulting atomic measure has moment vector \(b\), and contact at every atom makes its objective \((\lambda^-)^{\top}b=L_m(b)\). By (1), \(g_m(T_-)\) lies in the affine hyperplane \(\mathbf 1^{\top}x=1\) of dimension \(N-1\), so
\[
 r+1\leq N=N_m.
\]
The same reasoning with the upper contact set proves the upper endpoint assertions. Applying \(\text{lem:scm-moment-reduction}\) to these atomic measures produces the claimed compatible endpoint SCMs and completes the proof.